% ---- pg-558 （书页 546） ----

\noindent $K=4$

\begin{mvtable}
\footnotesize
\setlength{\tabcolsep}{4pt}%
\begin{tabular}{ccp{9.2cm}}
\toprule
聚类 & 聚类中的公司数目 & \multicolumn{1}{c}{公\quad\quad 司}\\
\midrule
1 & 5 & 爱达荷电力公司(8), 内华达电力公司(11), 普盖特桑德电力与电灯公司(16), 弗吉尼亚电气与电力公司(22), 肯塔基公用事业公司(9)\\
2 & 6 & 中央路易斯安那电气公司(3), 俄克拉何马天然气与电力公司(14), 南方公司(18), 得克萨斯公用事业公司(19), 亚利桑那公共服务公司(1), 佛罗里达电力与电灯公司(6)\\
3 & 5 & 新英格兰电气公司(12), 太平洋天然气与电力公司(15), 圣地亚哥天然气与电力公司(17), 联合照明公司(21), 夏威夷电气公司(7)\\
4 & 6 & 联合爱迪生公司(纽约)(5), 波士顿爱迪生公司(2), 麦迪逊天然气与电力公司(10), 北方电力公司(13), 威斯康星电力公司(20), 联邦爱迪生公司(4)\\
\bottomrule
\end{tabular}
\end{mvtable}

\begin{mvtable}
\mvtabcap{聚类中心之间的距离}
\[
\begin{array}{ccccc}
 & 1 & 2 & 3 & 4\\
1 & 0 & & & \\
2 & 3.08 & 0 & & \\
3 & 3.29 & 3.56 & 0 & \\
4 & 3.05 & 2.84 & 3.18 & 0
\end{array}
\]
\end{mvtable}

\noindent $K=5$

\begin{mvtable}
\footnotesize
\setlength{\tabcolsep}{4pt}%
\begin{tabular}{ccp{9.2cm}}
\toprule
聚类 & 聚类中的公司数目 & \multicolumn{1}{c}{公\quad\quad 司}\\
\midrule
1 & 5 & 内华达电力公司(11), 普盖特桑德电力与电灯公司(16), 爱达荷电力公司(8), 弗吉尼亚电气与电力公司(22), 肯塔基公用事业公司(9)\\
2 & 6 & 中央路易斯安那电气公司(3), 得克萨斯公用事业公司(19), 俄克拉何马天然气与电力公司(14), 南方公司(18), 亚利桑那公共服务公司(1), 佛罗里达电力与电灯公司(6)\\
3 & 5 & 新英格兰电气公司(12), 太平洋天然气与电力公司(15), 圣地亚哥天然气与电力公司(17), 联合照明公司(21), 夏威夷电气公司(7)\\
4 & 2 & 联合爱迪生公司(纽约)(5), 波士顿爱迪生公司(2)\\
5 & 4 & 联邦爱迪生公司(4), 麦迪逊天然气与电力公司(10), 北方电力公司(13), 威斯康星电力公司(20)\\
\bottomrule
\end{tabular}
\end{mvtable}

\begin{mvtable}
\mvtabcap{聚类中心之间的距离}
\[
\begin{array}{cccccc}
 & 1 & 2 & 3 & 4 & 5\\
1 & 0 & & & & \\
2 & 3.08 & 0 & & & \\
3 & 3.29 & 3.56 & 0 & & \\
4 & 3.63 & 3.46 & 2.63 & 0 & \\
5 & 3.18 & 2.99 & 3.81 & 2.89 & 0
\end{array}
\]
\end{mvtable}

$K=5$ 时的各种聚类轮廓图示于\hyperref[fig:12.12]{图 12.12}. 图中八个变量按组间变差与组内变差之比从大到小的顺序排列(对单变量的 $F$ 比, 见 6.4 节). 我们有
